Une fibre optique est formée d'un cœur en verre d'indice \num{\v{n1}} entouré d'une gaine d'indice \num[minimum-decimal-digits=2]{\fpeval{\v{n1} - \v{dn}}}. La lumière est guidée dans le cœur par réflexion totale interne. La vitesse de la lumière dans le vide est de \qty{3.00e8}{m/s}.

\begin{enumerate}[label=\alph*)]
    \item Déterminez l'angle critique à la frontière entre le cœur et la gaine. Pourquoi la gaine doit-elle avoir un indice plus petit que le cœur ?
    \item Déterminez la vitesse de la lumière dans le cœur.
    \item Combien de temps, au minimum, une impulsion lumineuse met-elle pour parcourir une fibre de \qty{\v{L}}{km} ?
\end{enumerate}
